% !TEX root = ../main.tex
% =====================================================================
% APPENDIX C -- REPRODUCIBILITY AND IMPLEMENTATION DETAILS
%
% The details that do not belong in the main methodology section: the
% full hyperparameter list, the engineering of the sweep, the exact
% commands and the two environments.
% =====================================================================

\section{Reproducibility and Implementation Details}
\label{app:repro}

\subsection{Full hyperparameters}
\label{app:hparams}

Every model in this report is one YAML file under \texttt{configs/models/}
(\Cref{sec:setup:pipeline}). \Cref{tab:hparams} lists the values shared by every configuration
trained, and separates the one group that is not: the encoder width and depth, which differ between
the wide base encoder and the smaller variants used in \Cref{sec:results}.

\begin{table}[htbp]
\caption{Hyperparameters shared by every model configuration, plus the two encoder sizes used.}
\label{tab:hparams}
\begin{center}
\small
\begin{tabular}{p{0.28\linewidth} p{0.30\linewidth} p{0.32\linewidth}}
\toprule
\textbf{Group} & \textbf{Parameter} & \textbf{Value} \\
\midrule
Encoder, wide (\texttt{sv2/seanet.yaml}) & width $\dmodel$ & 128 \\
        & number of blocks             & 6 \\
Encoder, narrow (sv4 family, used for every model in \Cref{sec:results}) & width $\dmodel$ & 64 \\
        & number of blocks             & 4 \\
\emph{Shared by both} & kernel sizes $\mathcal{K}$   & $\{5, 11, 23\}$, summed per block \\
        & dilation cap $\delta_{\max}$ & 16 \\
        & dropout                      & 0.2 \\
\midrule
Pooling & attention width $\dattn$     & 8 \\
        & dropout                      & 0.2 \\
        & positional encoding          & on \\
        & Top-$k$ fraction $\kappa$ (\Cref{sec:pool:topk})     & 0.1 \\
        & initial $\tau_0$ (\Cref{sec:pool:softmax}) & 1.0 \\
        & initial sharpness $\beta_0$ (\Cref{sec:pool:adaptive}) & 0.3 \\
        & initial blend $\lambda_0$, attention-max (\Cref{sec:pool:attnmax}) & 0.5 \\
        & initial blend $\lambda_0$, dual-stream (\Cref{sec:pool:dsmil}) & 0.05 \\
\midrule
Training& optimiser, learning rate     & Adam, $1.25\!\times\!10^{-3}$ \\
        & weight decay                 & $1.2\!\times\!10^{-4}$ \\
        & batch size                   & $\mathrm{clamp}(n_{\text{train}}/10, 2, 16)$ \\
        & epochs, patience             & 400, 60 \\
        & label smoothing              & 0.13 \\
        & focus penalty $\lambda_{\text{focus}}$ (\Cref{eq:loss}) & 0.01 \\
        & seeds                        & 3 for the four models in \Cref{tab:seeds}, 1 for the rest \\
\bottomrule
\end{tabular}
\end{center}
\end{table}
% Source: configs/models/sv2/seanet.yaml (wide encoder + training recipe), configs/models/sv4/*.yaml
% (narrow encoder, d=64 / n_blocks=4 -- shared by every model in Section 5, only the block type
% differs). Kernels, dilation cap, dropout and the whole training recipe are copy-pasted identically
% across sv1..sv7 on purpose, so results only change encoder-by-encoder or pooling-by-pooling.
% The four head-specific initial values are the constructor defaults in seanet/pooling.py:
% temperature=1.0 (line 234), init_beta=0.3 (284), init_lam=0.5 for AttentionMaxPooling (375) and
% init_lam=0.05 for DualStreamConjunctive (494). None is overridden in any YAML file.

\subsection{Running a study of this size}
\label{app:engineering}

Three engineering properties made the study of \Cref{sec:results} possible within the compute
constraints of \Cref{app:constraints}, and they affect what the stored results contain.

\begin{itemize}\itemsep2pt
  \item \textbf{Per-model result directories.} Each model writes one row per dataset per seed into
        its own directory, so two models can never overwrite each other and a partially finished
        model is still readable.
  \item \textbf{Resumable multi-dataset runs.} Grid'5000 jobs have a wall-clock limit and a
        complete run over the archive does not always finish within it, so each model keeps a record of the datasets it has
        already completed. Restarting the same job skips those, and an interrupted job costs at most
        the one dataset that was training when it stopped.
  \item \textbf{Tracking and curves.} Every run is logged to MLflow and saves its per-epoch loss
        curve next to its results, so the curve travels with the results when they are copied off
        the training machine.
\end{itemize}

\subsection{Commands}
\label{app:commands}

The complete study is reproduced by the commands below, run from the project root against one of the
configuration files under \texttt{configs/models/}. No step requires editing any code file.

{\footnotesize
\begin{verbatim}
# train one model over WebTraffic + the UCR archive (resumable)
python main.py train --model sv4/seanet_bottleneck_topk

# add a repeat with a different seed (results are stored per seed)
python main.py train --model sv4/seanet_bottleneck_topk --seed 1

# WebTraffic only, for a screening run or an ablation
python main.py webtraffic --model sv7/seanet_topk_k025

# cost numbers at the real WebTraffic shape
python scripts/profile_models.py --length 1008 --classes 10 --batch 32

# per-sample explanation figures and the page-1 teaser
python main.py interpret --model sv4/seanet_bottleneck_topk
python main.py teaser --models sv1/millet,sv4/seanet_bottleneck_topk

# seed spread, paired Wilcoxon test and the ensemble of Table 14
python scripts/ensemble_vote.py \
       --models sv4/seanet_bottleneck_topk \
                sv4/seanet_inputgate_adaptive \
       --baseline sv1/millet

# rebuild every derived file, in this order
python main.py leaderboard
python main.py results
python main.py report
python main.py paper
\end{verbatim}}

The order of the last four matters: \texttt{paper} reads the leaderboard, and the leaderboard reads
every model's \texttt{results.csv}. Every figure and table in this report is produced by
\texttt{python main.py paper} directly from those CSV files, so no number here was transcribed by
hand from a plot or a printed line.

\subsection{Software and hardware environment}
\label{app:env}

Development -- writing and testing the encoder and pooling code -- was done on a local Windows
machine with a 6\,GB RTX GPU, used for small tests only. The 72-model study ran on Grid'5000
\citep{balouek2013grid5000}, split across two sites so that a queue at one did not stall it.
\Cref{tab:environment} lists the two environments separately, because they are not identical and only
the Grid'5000 column produced the numbers in this report.

\begin{table}[htbp]
\caption{The two environments. Only the Grid'5000 column produced the results in this report.}
\label{tab:environment}
\begin{center}
\small
\begin{tabular}{p{0.20\linewidth} p{0.32\linewidth} p{0.38\linewidth}}
\toprule
 & \textbf{Local (development)} & \textbf{Grid'5000 (training)} \\
\midrule
OS              & Windows 11 & Linux \\
Python          & 3.10.20 & 3.10.20 \\
Conda env       & \texttt{seanet} & \texttt{seanet} (same on both sites) \\
PyTorch         & 2.0.1 & 2.0.1 \\
CUDA            & cu118 11.8 & cu118 11.8 \\
GPU             & 1$\times$ RTX, 6\,GB, small tests only & depends on the reserved node \\
Sites / clusters & -- & Lille (\texttt{chuc}, \texttt{chifflot}); Sophia
                        (\texttt{esterel22}, \texttt{esterel40}, \texttt{esterel43}) \\
Job scheduler   & --   & OAR \\
\bottomrule
\end{tabular}
\end{center}
\end{table}
% Source: requirements_versioned.txt (torch==2.0.1, mlflow==3.14.0, optuna==4.9.0, pandas==2.0.3,
% numpy==1.24.4, scikit-learn==1.3.0) and GRID5K_CMD_HELP.md (conda env "seanet"; -p chuc default
% Lille cluster with chifflot as an alternative; -p "cluster='esterel40'" / 'esterel43' for Sophia;
% OAR via oarsub). esterel22 added because the 2026-08-10 rebuild ran there.

\subsection{Code availability}
\label{app:code}

The codebase is organised as one module per stage of the pipeline under \texttt{seanet/} -- data
loading, the encoder and pooling implementations (\Cref{sec:methods}), training, and results and
figure generation -- driven by the configuration files under \texttt{configs/} (\Cref{app:hparams}).
It is released under the Apache License 2.0, and per its \texttt{NOTICE} file it is built on top of
Amazon's \texttt{MILTimeSeriesClassification} codebase (Copyright Amazon.com, Inc.\ or its
affiliates), the original implementation of \millet{} \citep{early2024millet}, so any redistribution
must preserve that attribution and its licence terms.

The repository is not yet public, but access can be granted to the examiner on request:
\url{https://github.com/ubaidur404786/sea-net}.
